\documentclass[11pt]{article}
\usepackage{amsmath,amssymb}
\usepackage[margin=0.85in]{geometry}

\title{From IHAT Redemption to Anchor-Bound Quality Reports}
\author{}
\date{}

\newcommand{\Ciss}{C_{\mathrm{iss}}}
\newcommand{\Crep}{C_{\mathrm{rep}}}
\newcommand{\Hash}{\mathsf{Hash}}
\newcommand{\Enc}{\mathsf{Encode}}

\begin{document}
\maketitle

\noindent\textbf{Status.} This is an experimental EC-IHAT/EC-ACT profile,
not a description of the existing version-one wire protocols. It describes
one successful Client--Moderator exchange with a report submitted at credential
presentation. The Moderator also acts as DAP Leader and Collector; a separate
Helper holds the other aggregation share. The proofs and report check described
below require distinct protocol identifiers and a specified DAP integration.
Message names below follow the MOLE and DAP drafts where applicable; any
combined upload or additional field is explicitly a \emph{proposed extension}.

\section{Task and notation}

The Moderator publishes an authenticated, ordered Anchor set
\begin{equation*}
  \mathcal A=(\mathsf{pk}_0,\ldots,\mathsf{pk}_{n-1}),
\end{equation*}
its epoch and credential context, and the report task configuration. The
experimental report profile uses two aggregators, Field128, and one Prio3
proof; a Field64 variant is not specified here. The
Client holds an endorsement
from Anchor $i$ and defines the \emph{private} tag $j=i+1$. The ordered set
and the meaning of $j$ remain fixed for the credential's lifetime. A new
set or mapping requires a new task and credential context.

All commitments below use the same prime-order group and fixed, independent
generators $H,R$. Let $B_{\mathrm I}$ be the EC-IHAT rerandomization
generator, and let $B_{\mathrm A},H_1,H_2,H_3,H_4$ be the EC-ACT generators.
$H_j$ is a new ACT generator for the hidden Anchor attribute. The Moderator
holds ACT signing secret $x_M$. Distinct domain separators and canonical,
versioned encodings are required throughout. The Client draws fresh and
independent blinding factors $\rho_{\mathrm{iss}}$ and $\rho_{\mathrm{rep}}$;
it never reuses either commitment.

\section{Message map: MOLE, DAP, and the link}

Each entry names a protocol message, its MOLE and DAP portions, and the
cross-protocol check. ``None'' means that phase has no DAP report. The
origin receipt is application data, not an existing MOLE or DAP message.
Helper exchanges may be batched into DAP aggregation jobs.

\begin{description}
\item[M1: Client--Anchor.] \textbf{MOLE:} Endorsement grant and response.
  \textbf{DAP:} None. \textbf{Link:} The issuing Anchor is
  $\mathsf{pk}_i$ in the configured set; no report exists yet.
\item[M2: Moderator--Client.] \textbf{MOLE:} \texttt{RedemptionChallenge}.
  \textbf{DAP:} None. \textbf{Link:} Select the accepted Anchor set,
  epoch, credential context, and binding profile.
\item[M3: Client--Moderator.] \textbf{MOLE:} \texttt{Redemption} and
  \texttt{IssuanceRequest} in one \texttt{CredentialRequest}.
  \textbf{DAP:} None. \textbf{Link:} Both payloads carry identical
  $\Ciss$; the IHAT proof binds the issuing Anchor to that commitment and
  the exact ACT request digest. The ACT proof binds $\Ciss$ to hidden $j$
  in $K$.
\item[M4: Moderator--Client.] \textbf{MOLE:} \texttt{IssuanceResponse}.
  \textbf{DAP:} None. \textbf{Link:} The Moderator signs the $K$ accepted
  in M3; the Client verifies against its saved $K$ and retains $j$.
\item[M5: Moderator--Client.] \textbf{MOLE:} Credential \texttt{Challenge}
  (extended for this profile). \textbf{DAP:} Task configuration, not a
  report. \textbf{Link:} Fix the spend policy, DAP task, origin, and
  label-policy scope of the presentation.
\item[M6: Origin--Client.] \textbf{MOLE:} None (application outcome receipt).
  \textbf{DAP:} None. \textbf{Link:} Authenticate the origin's label and
  policy version so the Client cannot choose or substitute them.
\item[M7: Client--Moderator.] \textbf{MOLE:}
  \texttt{PresentationAndUpdate} with spend proof.
  \textbf{DAP:} \texttt{Report} in \texttt{UploadRequest}, in the same HTTP
  submission (proposed combined transport). \textbf{Link:} The MOLE
  proof binds hidden $j$ to $\Crep$ and the digest of the exact DAP report
  and outcome scope. A public DAP extension carries $\Crep$; each encrypted
  private extension carries its $\rho_m$.
\item[M8: Moderator--Helper.] \textbf{MOLE:} Moderator accepts M7's
  presentation and records the credential result.
  \textbf{DAP:} \texttt{AggregationJobInitReq} with \texttt{VerifyInit}
  (extended). \textbf{Link:} Send the encrypted Helper
  \texttt{ReportShare}, the Leader's Prio3 verifier share and $P_0$, and
  an authenticated Moderator acceptance assertion bound to the task,
  report ID, report digest, and $\Crep$. The Helper cannot independently
  verify the ACT signing-key proof.
\item[M9: Helper--Moderator.] \textbf{MOLE:} None.
  \textbf{DAP:} Aggregation-job response (new versioned encoding).
  \textbf{Link:} The Helper computes the verifier message and its own
  $P_1$, checks Prio3 and $P_0+P_1\in\{\Crep,\Crep+pH\}$, and returns
  the verifier message and $P_1$. The Leader performs the same checks
  before counting its output share. No extra preparation round trip is
  required for one-round Prio3.
\item[M10: Moderator--Client.] \textbf{MOLE:} \texttt{Update} or recorded
  \texttt{NoUpdate}. \textbf{DAP:} No report message.
  \textbf{Link:} Any replacement credential retains hidden $j$; define
  atomic result recording, rejected reports, and byte-identical retries.
  No client-visible report authorization token is sent.
\item[M11: Moderator--Helper.] \textbf{MOLE:} None (Moderator is also
  Collector). \textbf{DAP:} \texttt{AggregateShareReq} for a later batch.
  \textbf{Link:} Both sides agree on the accepted report set and scoped
  task; DAP count, checksum, and batch-release checks apply.
\item[M12: Helper--Moderator.] \textbf{MOLE:} None.
  \textbf{DAP:} \texttt{AggregateShare}.
  \textbf{Link:} Helper encrypts to the Moderator as Collector, which
  combines it with its own share and learns the
  released per-Anchor histogram for this origin and label.
\end{description}

M7 is \emph{not} an existing MOLE HTTP message that automatically uploads a
DAP report. The combined transport, public commitment extension, private
blinding-share extensions, mandatory acceptance assertion in M8, and
commitment points piggybacked on M8--M9 must be specified. DAP has an
aggregation-job request extension but no matching response extension;
M9 needs a new versioned response encoding, without changing the VDAF
verifier message. Ordinary
\texttt{Prio3Histogram} and unextended DAP do not verify these bindings.

\section{IHAT redemption and credential request}

Messages M2--M3 are MOLE payloads; they contain no DAP \texttt{Report}.
The Client prepares the ACT request first so that its exact encoding can be
bound into the IHAT proof, then sends both payloads in M3.

\paragraph{Client, before sending the combined request.}
The Client chooses $\rho_{\mathrm{iss}}$ and constructs
\begin{equation}
  \Ciss=jH+\rho_{\mathrm{iss}}R.
\end{equation}
It prepares the ACT issuance commitment and a proof $\pi_{\mathrm{issue}}$
of knowledge of the \emph{same} $j$ in both equations:
\begin{equation}
  K=kH_2+rH_3+jH_j,
  \qquad
  \Ciss=jH+\rho_{\mathrm{iss}}R.
  \label{eq:issue}
\end{equation}
The versioned ACT issuance request contains $(K,\Ciss,\pi_{\mathrm{issue}})$.
The Client computes its digest $d_{\mathrm{issue}}$ over the complete encoded
request, so that IHAT redemption can refer to these exact request bytes.

The Client rerandomizes the endorsed IHAT key to $\widehat X$ and supplies
the ordinary endorsement redemption proof together with a new, joint
same-Anchor proof. The latter proves the following statement without
revealing $i$, $\delta$, or $\rho_{\mathrm{iss}}$:
\begin{equation}
  \exists\,i,\delta,\rho_{\mathrm{iss}}:\quad
  0\leq i<n,\qquad
  \widehat X-\mathsf{pk}_i=\delta B_{\mathrm I},\qquad
  \Ciss-(i+1)H=\rho_{\mathrm{iss}}R.
  \label{eq:ihat}
\end{equation}
Both equations use the \emph{same branch} $i$. An OR proof over the set is
one possible realization. Its transcript binds the ordered set and epoch,
IHAT contexts and challenge, the endorsement redemption, $\Ciss$, and
$d_{\mathrm{issue}}$. Adding only $\Ciss$ to an ordinary IHAT transcript
would not prove Equation~\eqref{eq:ihat}.

\paragraph{Moderator.}
On receipt of M3, the Moderator verifies the ordinary IHAT redemption, the joint proof, the
challenge and context, the exact ACT request digest, and
$\pi_{\mathrm{issue}}$. It checks that both requests carry the identical
$\Ciss$ and atomically checks and records the IHAT nullifier in the
appropriate replay scope. It learns that an endorsed Anchor has a valid
hidden index $j$, but neither its value nor the Anchor's identity. The joint IHAT proof
binds $\Ciss$ to that Anchor; Equation~\eqref{eq:issue} binds the same
commitment to the credential attribute.

\section{Credential issuance}

\paragraph{Moderator.}
After accepting M3, the Moderator chooses the initial
balance $c$ under its policy and signs the ACT message
\begin{equation}
  X_A=B_{\mathrm A}+cH_1+\mathsf{ctx}_{\mathrm{cred}}H_4+K.
  \label{eq:signed}
\end{equation}
For the illustrative EC-ACT signature relation, its response contains
$A=(x_M+e)^{-1}X_A$ and the associated signing proof and parameters. The
Moderator does not need to open $K$ or learn $j$ to sign it. This is
MOLE message M4; it carries no DAP payload.

\paragraph{Client.}
The Client verifies the response against the exact saved $K$ and credential
context, then stores $(k,c,r,j,A,e)$. Because $K$ in
Equation~\eqref{eq:signed} contains $jH_j$, a later proof of possession
of this credential can authenticate $j$ without disclosing it. The
credential issuance proof and signature, not public equality of
commitments, carry the binding forward.

\section{Credential presentation and report submission}

An origin supplies an authenticated outcome label and policy version for an
eligible interaction. This section assumes the outcome is known at
presentation time. The task fixes the origin and label scope; reports with
incompatible label semantics are not combined. The Moderator's MOLE
challenge M5 selects the policy, and the origin supplies application
receipt M6. Neither is itself a DAP report.

\paragraph{What binds the report to the presentation.}
The Client submits its ACT presentation, authenticated origin receipt, and
the exact DAP \texttt{Report} together. The ACT proof uses one hidden witness
$j$ in both the signed credential and the fresh public report commitment
$\Crep=jH+\rho_{\mathrm{rep}}R$. Its transcript binds the task, report ID,
origin and label policy, and a digest of the whole encoded DAP report,
including both encrypted shares. The Moderator checks that the commitment
proved in ACT is the one in the report's proposed public extension. This
identifies the \emph{particular} report authorized by the credential, but
cannot by itself show which bucket its encrypted VDAF shares count. The
aggregators supply that final link using the point check in M8--M9; neither
point is supplied by the Client in the DAP report. This combined submission
is a proposed transport profile, not an existing DAP upload.

\paragraph{Client, preparing one submission.}
The Client draws a new $\rho_{\mathrm{rep}}$ and computes
\begin{equation}
  \Crep=jH+\rho_{\mathrm{rep}}R.
  \label{eq:report}
\end{equation}
For its private histogram bucket $i=j-1$, it forms a one-hot vector
$v=(v_0,\ldots,v_{n-1})$ with $v_i=1$.
It creates two ordinary Prio3Histogram input shares of $v$ and private
shares of $\rho_{\mathrm{rep}}$, encrypting one pair for the Moderator and the other
for the Helper. In DAP terms, the report's metadata includes a fresh
\texttt{report\_id}, which is also the VDAF nonce, and a public extension
carrying $\Crep$. The DAP \texttt{Report} includes the VDAF public share
and two HPKE-encrypted \texttt{PlaintextInputShare} values; each plaintext
contains that aggregator's VDAF input share and a proposed private extension
carrying $\rho_m$. The task ID identifies the DAP task and versioned
Anchor mapping. The Client computes a domain-separated digest
\begin{align}
  d_{\mathrm{rep}}={}&\Hash(\mathsf{task\_id}\,\|\,\mathsf{origin}\,\|\,
    \mathsf{label\_policy}\,\|\,\nonumber\\
    &\qquad\Enc(\mathsf{outcome\_receipt})\,\|\,
    \Enc(\mathsf{DAP\ Report})).
  \label{eq:digest}
\end{align}
This digest covers the exact report ID, $\Crep$, public share, and both
ciphertexts through the canonical DAP \texttt{Report} encoding. It excludes
the ACT spend proof itself to avoid a circular hash dependency.

The Client proves possession and valid use of the credential in an ACT
presentation whose transcript binds $d_{\mathrm{rep}}$, the task, report
ID, origin, and label-policy version. Schematically, the spend proof
demonstrates that the old credential signature authenticates the hidden
$j$ and that \emph{the same witness} satisfies
\begin{align}
  K_{\mathrm{next}}&=k_{\mathrm{next}}H_2+r_{\mathrm{next}}H_3+jH_j,
    \label{eq:next}\\
  \Crep&=jH+\rho_{\mathrm{rep}}R.
    \label{eq:same-j}
\end{align}
Here $K_{\mathrm{next}}$ denotes the hidden-attribute portion of the
replacement credential commitment; the ACT proof also contains its usual
signature-possession, balance, spend, and allowance constraints. In the
experimental EC-ACT relation, the signature and hidden-attribute checks
can be written more concretely as follows. The Client forms
$A'=r_1r_2 A$ and
$\bar B=r_1(B_{\mathrm A}+cH_1+kH_2+rH_3+
\mathsf{ctx}_{\mathrm{cred}}H_4+jH_j)$; the Moderator computes
$\bar A=x_M A'$ and
$H'_1=B_{\mathrm A}+kH_2+\mathsf{ctx}_{\mathrm{cred}}H_4$.
The proof establishes
\begin{align}
  \bar A&=x_M A'=-eA'+r_2\bar B,\\
  H'_1&=r_3\bar B-cH_1-rH_3-jH_j,
    \qquad r_3=r_1^{-1},\\
  K_n&=k_{\mathrm{next}}H_2+r_nH_3+jH_j,\\
  \Crep&=jH+\rho_{\mathrm{rep}}R.
  \label{eq:spend}
\end{align}
In M7, the Client submits the MOLE \texttt{PresentationAndUpdate}, origin
receipt, and the DAP \texttt{Report} \emph{together} to the Moderator.
The same $\Crep$ appears in the MOLE proof and the DAP public extension;
the MOLE transcript binds Equation~\eqref{eq:digest}. This is a new
combined-transport profile, not an existing DAP upload authenticated by an
ordinary MOLE header. No client-visible intermediate report-authorization
token is needed.

\paragraph{Moderator.}
For M7, the Moderator verifies the credential presentation under the correct signing
key and context; checks the authenticated outcome label; checks that the
commitment in the proof is the $\Crep$ in the submitted payload; recomputes
$d_{\mathrm{rep}}$; and enforces the credential-nullifier and report-ID
replay rules. It learns neither $j$ nor the individual histogram bucket.
It then acts as DAP Leader in M8, forwarding the Helper's encrypted
\texttt{ReportShare} and authenticated common metadata, including a
Moderator assertion tied to the accepted MOLE presentation and the exact
task, report ID, $\Crep$, and $d_{\mathrm{rep}}$. The Helper relies on the Moderator's
authenticated acceptance decision for the ACT proof, but independently
processes its own VDAF share. For one-round Prio3, the Leader sends its
verifier share and $P_0$ in M8; the Helper verifies, computes the verifier
message and $P_1$, and returns both in M9. The Leader verifies on receipt.
These are extensions to the existing request and response, not a second
round trip. Both aggregators run the
Prio3 verification protocol on
the full shares to establish validity
\begin{equation}
  v_u\in\{0,1\}\ (0\leq u<n),\qquad
  \sum_{u=0}^{n-1}v_u=1,
  \label{eq:vdaf}
\end{equation}
and the prescribed binding of their projected index and blinding shares to
$\Crep$ before either counts an output share. Concretely, if the VDAF field
has modulus $p$, the group has order $q>2p$, and the actual private
measurement shares are $v_{m,u}$, each aggregator derives locally
\begin{equation}
  t_m=\sum_{u=0}^{n-1}(u+1)v_{m,u}\pmod p.
\end{equation}
Linearity and Equation~\eqref{eq:vdaf} imply
$t_0+t_1=j\pmod p$. Each aggregator then computes
\begin{equation}
  P_m=\mathsf{int}(t_m)H+\rho_mR,
  \qquad \rho_0+\rho_1=\rho_{\mathrm{rep}}\pmod q.
\end{equation}
Both check that $P_0+P_1$ is either $\Crep$ or $\Crep+pH$: the latter
case accounts for a carry when canonical field shares sum to $t+p$.
This point exchange is an additional, mandatory aggregation-protocol step,
not part of ordinary Prio3Histogram verification. No auxiliary VDAF tag,
new validity circuit, or new VDAF algorithm ID is needed. Only the $v$
coordinates enter the histogram aggregate. The Moderator, acting as
Collector, combines aggregate shares to learn counts by Anchor for this
origin and outcome label, subject to the task's release policy.

In MOLE response M10, if the Moderator refunds or updates the credential,
its response signs a
replacement commitment containing the \emph{same hidden $j$} as
Equation~\eqref{eq:next}. The Client checks the response and retains $j$
with the updated balance and secrets. A later report uses another fresh
commitment as in Equation~\eqref{eq:report}. This response need not wait
for batch collection, but the profile must say how verification failure,
result recording, and byte-identical presentation retries interact with
report admission. Later, in DAP messages M11--M12, the Moderator as
Collector requests the Helper's encrypted aggregate share for an eligible
batch and combines it with its own share.

\section{What actually links the two commitments?}

No party compares $\Ciss$ and $\Crep$ directly: their independent
blindings deliberately make them unlinkable as group elements. Instead,
Equation~\eqref{eq:ihat} proves that $\Ciss$ contains the redeemed
Anchor's tag; Equation~\eqref{eq:issue} places that tag in the signed
credential; and Equations~\eqref{eq:spend} and~\eqref{eq:same-j} prove
that the credential's tag is the one in $\Crep$. Equation~\eqref{eq:vdaf},
the linear projection, and the aggregation binding check finally tie it to
the counted bucket.
Each of these links is necessary. A digest or transcript reference alone
does not prove equality of the hidden tags.

The commitment check enforces report integrity against a cheating Client
only if both aggregators execute it on their actual shares; a corrupt
aggregator can forge its check point. A non-colluding Helper is required
for individual-report privacy. The effect of revealing which field-carry
comparison succeeds also needs analysis. These are limitations of the
illustrative construction, not properties supplied by ordinary DAP.

For reporting \emph{only at IHAT redemption}, credential issuance and
presentation are unnecessary for this purpose. A redemption proof could
instead bind the actual secret-shared report directly to the IHAT selector,
as in the experimental VOLE-IHAT share-reporting approach. This note covers
the different case of carrying the Anchor into a later credential report.

\end{document}
